\documentclass[11pt,a4paper]{article}
\usepackage[T1]{fontenc}
\usepackage[utf8]{inputenc}
\usepackage{amsmath,amssymb,amsthm}
\usepackage[margin=1in]{geometry}
\usepackage[colorlinks=true,linkcolor=blue,urlcolor=blue]{hyperref}
\theoremstyle{plain}
\newtheorem{theorem}{Theorem}
\newtheorem{lemma}[theorem]{Lemma}
\newtheorem{proposition}[theorem]{Proposition}
\newtheorem{corollary}[theorem]{Corollary}
\theoremstyle{definition}
\newtheorem{remark}[theorem]{Remark}

\title{A condition on the multiset of every subblock: OEIS A184203}
\author{Adrian Perez Fontelles\\ \small Independent researcher}
\date{7 September 2026}
\begin{document}
\maketitle

\begin{abstract}
OEIS A184203 carries an empirical recurrence of order $36$ contributed by R.~H.~Hardin. It is
true, and it is decidable rather than empirical. The entry's condition on a subblock depends
only on which values the block holds and how many times, not on where they sit, so it is
decided by the few consecutive lines the block spans. The admissible windows are the vertices
of a finite digraph and the arrays counted are exactly its walks; hence $a(n)$ is a walk count
on $S=36$ vertices and is $C$-finite, and whether the conjectured recurrence annihilates it
is settled by a finite exact computation, with the Cayley--Hamilton theorem bounding the
work.
\end{abstract}

\noindent\small 2020 Mathematics Subject Classification. 05A15, 05B45, 15B36.\normalsize

\section{The conjecture}

OEIS A184203 is ``Number of (n+1) X 6 binary arrays with no 2 X 2 subblock containing fewer than two 1's.''. It has offset $1$ and begins
\[
1049, 29593, 721939, 18588589, 471487297, 12021784217, 306055720601, 7795714368919,\ \dots
\]
The entry states, as a conjecture contributed by its author and never marked settled:
\begin{quote}
Empirical: a(n) = 26*a(n-1)+103*a(n-2)-3175*a(n-3)+843*a(n-4)+124668*a(n-5)-241004*a(n-6)-1741528*a(n-7)+4512428*a(n-8)+11666399*a(n-9)-36429185*a(n-10)-41916935*a(n-11)+161955426*a(n-12)+79119202*a(n-13)-435611222*a(n-14)-51541960*a(n-15)+738522868*a(n-16)-83869087*a(n-17)-800207178*a(n-18)+216307031*a(n-19)+552396475*a(n-20)-212199576*a(n-21)-239220431*a(n-22)+114209818*a(n-23)+63085385*a(n-24)-36259721*a(n-25)-9488703*a(n-26)+6873528*a(n-27)+671231*a(n-28)-762566*a(n-29)+22*a(n-30)+47026*a(n-31)-2787*a(n-32)-1453*a(n-33)+139*a(n-34)+17*a(n-35)-2*a(n-36).
\end{quote}
As of the ``Last modified'' line on the live entry (May 08 2022 22:55:20, revision 7) this is still
recorded as empirical, and nothing on the entry records it as proved.

\section{The count is a walk count}

The entry counts arrays with $6$ columns over the alphabet $\{0,1\}$, growing one row at a time, and asks that every $2\times2$ subblock holds at least $2$ entries equal to $1$.

That condition does not depend on where in the block a value sits, only on which values the block holds and how many times --- it is a function of the block's multiset of entries. In particular it is decided by $2$ consecutive rows, the ones the block spans, so carrying the previous $1$ row as the state is enough: each new row completes every block that ends on it, and no judgement is left over at the end of the array.

Merging vertices with identical futures leaves $S=36$. An admissible array is exactly a walk, so with $M$ the adjacency matrix and $\iota,\tau$ the vectors recording which states may begin and end an array,
\[
a(n)\;=\;\iota^{\!\top}M^{\,n}\tau ,
\]
the exponent fixed by the entry's own shape and checked against its published terms. In
particular $a$ is $C$-finite.

\section{The proof}

\begin{lemma}\label{lem:crit}
Let $q(t)=t^{36}-\sum_i c_it^{36-i}$ be the characteristic polynomial of the
conjectured recurrence. The residual $a(n)-\sum_ic_ia(n-i)$ equals
$\iota^{\!\top}M^{\,j}q(M)\tau$ for the corresponding walk index $j$, so the recurrence
holds from some point on if and only if $\iota^{\!\top}M^{j}q(M)\tau=0$ for all large $j$,
and it is enough to examine $j<S$.
\end{lemma}

\begin{proof}
Factor $M^{j}$ out of $M^{j+36}-\sum_ic_iM^{j+36-i}$. For the bound, $M^{S}$
is an integer combination of $I,M,\dots,M^{S-1}$ by Cayley--Hamilton, so the numbers
$u_j=\iota^{\!\top}M^{j}q(M)\tau$ obey the monic recurrence given by the characteristic
polynomial of $M$; $S$ consecutive zeros force every later one, and the last nonzero $u_j$
pins the threshold exactly. A constant factor in front of $a$ divides out of the residual and
changes nothing here.
\end{proof}

\begin{theorem}
\[
a(n)\;=\;26\,a(n-1) + 103\,a(n-2) - 3175\,a(n-3) + 843\,a(n-4) + 124668\,a(n-5) - 241004\,a(n-6) - 1741528\,a(n-7) + 4512428\,a(n-8) + 11666399\,a(n-9) - 36429185\,a(n-10) - 41916935\,a(n-11) + 161955426\,a(n-12) + 79119202\,a(n-13) - 435611222\,a(n-14) - 51541960\,a(n-15) + 738522868\,a(n-16) - 83869087\,a(n-17) - 800207178\,a(n-18) + 216307031\,a(n-19) + 552396475\,a(n-20) - 212199576\,a(n-21) - 239220431\,a(n-22) + 114209818\,a(n-23) + 63085385\,a(n-24) - 36259721\,a(n-25) - 9488703\,a(n-26) + 6873528\,a(n-27) + 671231\,a(n-28) - 762566\,a(n-29) + 22\,a(n-30) + 47026\,a(n-31) - 2787\,a(n-32) - 1453\,a(n-33) + 139\,a(n-34) + 17\,a(n-35) - 2\,a(n-36)
\]
for every $n>35$.
\end{theorem}

\begin{proof}
The vector $w=q(M)\tau$ was formed by $36$ matrix--vector products in exact integer
arithmetic and $\iota^{\!\top}M^{j}w$ evaluated for $j=0,1,\dots$, again exactly; those
integers vanish from the index corresponding to $n=35+1$ onwards and the run continued
until the working vector was identically zero, which settles every larger $j$ at once.
\end{proof}

\section{Verification}

Four checks, each independent of the algebra above.

First, the model reproduces all $14$ terms the entry publishes, exactly, in integer
arithmetic. This is what ties the model to the entry: the digraph is built by reading the
entry's English, and a misreading gives different counts.

Second, the reading was pinned before the digraph was written, by a program that enumerates
every array of the given shape one by one, forms each subblock directly and tests the entry's
condition on it. That program shares no code with the transfer matrix, so an error in one does
not hide an error in the other.

Third, the conjectured recurrence was evaluated directly on the published terms, with no
matrices involved, and holds wherever the proved range and the data overlap.

Fourth, the annihilation test of Lemma~\ref{lem:crit} was carried out in exact integer
arithmetic on the whole vector rather than on a sampled prefix.

\begin{thebibliography}{9}
\bibitem{oeis} The OEIS Foundation, \emph{The On-Line Encyclopedia of Integer
Sequences}, \url{https://oeis.org}, 2026. Sequence A184203.
\bibitem{stanley} R.~P.~Stanley, \emph{Enumerative Combinatorics}, Volume 1, 2nd ed.,
Cambridge University Press, 2011. (Transfer-matrix method, Section 4.7.)
\bibitem{fs} P.~Flajolet and R.~Sedgewick, \emph{Analytic Combinatorics}, Cambridge
University Press, 2009.
\bibitem{horn} R.~A.~Horn and C.~R.~Johnson, \emph{Matrix Analysis}, 2nd ed., Cambridge
University Press, 2013.
\end{thebibliography}

\end{document}
